\documentclass[11pt]{ctexart}
%=====================================================================
% theme-violet.tex —— 讲义精讲 · 紫花主题
% 用法：\documentclass[11pt]{ctexart}
%       \input{assets/theme-violet}
%       \begin{document} ... \end{document}
% 只用 ctexart 自带能力 + xcolor/enumitem/fancyhdr/pifont/amsmath，
% 不依赖 tikz/tcolorbox，保证 MiKTeX 可直接编译。
%=====================================================================

\usepackage[a4paper,left=2.0cm,right=2.0cm,top=2.1cm,bottom=2.3cm,%
            headsep=12pt,footskip=22pt]{geometry}
\usepackage{amsmath,amssymb}
\usepackage{xcolor}
\usepackage{array}
\usepackage{enumitem}
\usepackage{pifont}
\usepackage{fancyhdr}
\usepackage[hidelinks]{hyperref}

%---------------------------------------------------------------- 配色
\definecolor{violet}{RGB}{107,78,143}      % 主色：紫
\definecolor{violetdeep}{RGB}{63,44,92}    % 深紫：标题/强调
\definecolor{petal}{RGB}{183,159,208}      % 花瓣：装饰线、花饰
\definecolor{bandviolet}{RGB}{247,243,251} % 浅紫底：档案/提示框
\definecolor{bandcream}{RGB}{252,249,242}  % 米白底：易错/注意框
\definecolor{rulepurple}{RGB}{217,207,230} % 分隔线
\definecolor{inkgray}{RGB}{92,96,110}      % 次要文字

%---------------------------------------------------------------- 字体
% 有 Noto CJK 用 Noto，没有就退回 ctex 默认（Windows 下为宋体/黑体）
\IfFontExistsTF{Noto Serif CJK SC}{\setCJKmainfont{Noto Serif CJK SC}}{}
\IfFontExistsTF{Noto Sans CJK SC}{\setCJKsansfont{Noto Sans CJK SC}}{}

%---------------------------------------------------------------- 版式
\setlist{nosep,leftmargin=1.9em}
\linespread{1.28}
\setlength{\parskip}{3.5pt}
\setlength{\parindent}{0pt}
\setlength{\tabcolsep}{5pt}
\renewcommand{\arraystretch}{1.28}

%---------------------------------------------------------------- 花饰
% \ding{95} = 五瓣小花，\ding{96} = 星形花；都来自 pifont，必定存在
\newcommand{\orn}{\textcolor{petal}{\ding{95}}}          % 行内小标
\newcommand{\ornstar}{\textcolor{petal}{\ding{96}}}      % 星形花
\newcommand{\flowerrule}{%                                % 居中花饰分隔线
  \par\vspace{7pt}%
  \noindent\makebox[\linewidth]{\color{petal}%
    \rule{0.20\linewidth}{0.6pt}\hspace{8pt}{\small\ding{95}}\hspace{4pt}%
    \raisebox{1.1pt}{\tiny\ding{95}}\hspace{4pt}{\small\ding{95}}\hspace{8pt}%
    \rule{0.20\linewidth}{0.6pt}}%
  \par\vspace{7pt}}

%---------------------------------------------------------------- 标题与节头
\newcommand{\doctitle}[3]{%   #1 主标题  #2 副标题  #3 元信息（可用 \\ 分行）
  \begin{center}
    {\color{petal}\small \orn\enspace #2\enspace \ornstar}\\[7pt]
    {\LARGE\bfseries\color{violetdeep} #1}\\[10pt]
    {\color{petal}\rule{0.55\linewidth}{1pt}}\\[4pt]
    {\small\color{inkgray} #3}
  \end{center}\vspace{2pt}}

\newcommand{\hdr}[2]{%   #1 章节号（如 §0）  #2 标题
  \par\vspace{12pt}%
  \noindent\colorbox{violetdeep}{\color{white}\bfseries\small\;#1\;}\hspace{7pt}%
  {\color{violetdeep}\large\bfseries #2}\par\nobreak
  \vspace{3pt}{\color{petal}\hrule height 1.2pt}\vspace{6pt}}

\newcommand{\subhdr}[1]{%
  \par\vspace{6pt}%
  \noindent{\color{violet}\bfseries #1}\par\vspace{1pt}}

%---------------------------------------------------------------- 提示框
\newcommand{\callout}[1]{% 浅紫底：讲义档案、读法说明
  \par\medskip\noindent
  \fcolorbox{rulepurple}{bandviolet}{\parbox{\dimexpr\linewidth-2\fboxsep-2\fboxrule}{\vspace{2pt}#1\vspace{2pt}}}%
  \par\medskip}

\newcommand{\warnbox}[1]{% 米白底：易错、边界条件
  \par\medskip\noindent
  \fcolorbox{rulepurple}{bandcream}{\parbox{\dimexpr\linewidth-2\fboxsep-2\fboxrule}{\vspace{2pt}#1\vspace{2pt}}}%
  \par\medskip}

\newcommand{\gapitem}[2]{% #1 缺口类型（紫） #2 内容
  \par\vspace{2pt}\noindent
  \hangindent=1.4em\hangafter=1
  {\color{petal}\small\ding{96}}\enspace{\bfseries\color{violetdeep}#1}\enspace #2\par}

% 前提旁注：只放在「需要理解的那一步」后面，给陈述 + 本题代入，不重证。
\newcommand{\aside}[1]{%
  \par\vspace{1pt}\noindent
  \hangindent=1.6em\hangafter=1
  {\small\color{inkgray}%
    \hspace*{0.6em}{\color{petal}$\triangleright$}\hspace{2pt}#1}%
  \par\vspace{2pt}}

%---------------------------------------------------------------- 表格
% 本机 MiKTeX 的 colortbl 与 array 版本不匹配（\rowcolor 会报错），
% 故表格统一用「上粗紫线 + 表头 + 中细线」的三线式，不依赖 colortbl。
\newcommand{\tabtop}{\par\vspace{3pt}{\color{violet}\hrule height 1.2pt}\vspace{1pt}}
\newcommand{\tabmid}{\par\vspace{1pt}{\color{petal}\hrule height 0.5pt}\vspace{1pt}}
\newcommand{\tabbot}{\par\vspace{3pt}{\color{rulepurple}\hrule height 0.5pt}\vspace{8pt}}
\newcommand{\thead}[1]{{\bfseries\color{violetdeep}#1}}

%---------------------------------------------------------------- 页眉页脚
\pagestyle{fancy}
\fancyhf{}
\renewcommand{\headrulewidth}{0.4pt}
\renewcommand{\headrule}{\hbox to\headwidth{\color{rulepurple}\leaders\hrule height \headrulewidth\hfill}}
\fancyhead[L]{\small\color{inkgray}\songti \doctag}
\fancyhead[R]{\small\color{inkgray}\songti \coursetag}
\fancyfoot[C]{\small\color{inkgray}\orn\enspace{\color{violet}\thepage}\enspace\orn}

% 每篇讲义在导言区设这两个值
\newcommand{\doctag}{}
\newcommand{\coursetag}{}
\newcommand{\settags}[2]{\renewcommand{\doctag}{#1}\renewcommand{\coursetag}{#2}}

%---------------------------------------------------------------- 数学记号
\newcommand{\vect}[1]{\boldsymbol{#1}}
\newcommand{\dd}{\,\mathrm{d}}
\newcommand{\ddx}{\frac{\mathrm{d}}{\mathrm{d}x}}
\addtolength{\jot}{2pt}



\settags{Problem 1 · 常点幂级数}{ODE · Lecture 18--20}

% 本机 TeX 对 \dimexpr 组合长度的解析不稳定，盒子宽度改用固定比例
\renewcommand{\callout}[1]{\par\medskip\noindent
  \colorbox{bandviolet}{\parbox{0.955\linewidth}{\vspace{2pt}#1\vspace{2pt}}}%
  \par\medskip}
\renewcommand{\warnbox}[1]{\par\medskip\noindent
  \colorbox{bandcream}{\parbox{0.955\linewidth}{\vspace{2pt}#1\vspace{2pt}}}%
  \par\medskip}
\renewcommand{\qstem}[1]{\par\medskip\noindent
  \colorbox{bandviolet}{\parbox{0.955\linewidth}{\vspace{2pt}%
    {\small\bfseries\color{violet}题干}\hspace{0.8em}#1\vspace{2pt}}}%
  \par\medskip}

\setlength{\headheight}{22pt}
\begin{document}

\doctitle{常点处的幂级数解}{Problem 1 完整推导}{题面来源：Lecture 18--20 作业截图（截止：第九周周日 23:59）\quad·\quad 2026-10-28}

\qstem{在 $x=0$ 附近，用两个线性无关的幂级数解构造方程的通解；每个解写出前五个非零项，并指出级数保证收敛的区间。
\[ y''+(1-x)y'+(x^2-1)y=0 \]}

\hdr{§1}{这是常点还是奇点}

先把方程写成首项系数为 $1$ 的标准形状
\[ y''+P(x)y'+Q(x)y=0,\qquad P(x)=1-x,\quad Q(x)=x^2-1 . \]

\callout{
\textbf{定义（常点／奇点）}\;
若 $P,Q$ 在 $x_0$ 的某个邻域内都解析，称 $x_0$ 为该方程的\emph{常点}；只要 $P$ 或 $Q$ 在 $x_0$ 处不解析（有极点），$x_0$ 就是\emph{奇点}。\\
判据落在\textbf{系数函数}上，与“某个系数在 $x_0$ 取不取零”无关。
}

本题 $P,Q$ 都是多项式，在全平面解析，并且
\[ P(0)=1\neq 0,\qquad Q(0)=-1 . \]
所以 $x=0$ 是\textbf{常点}，不是奇点。常点处可以设普通幂级数
\[ y=\sum_{n=0}^{\infty}a_nx^n , \]
从 $a_0,a_1$ 出发一路递推下去。

\warnbox{
\textbf{两个容易混的点。}
（1）$x^2-1=0$ 的根 $x=\pm1$ 只是“$y$ 前面的系数取零”，系数函数本身仍解析，\textbf{不构成奇点}；它们不影响幂级数的收敛半径。\\
（2）什么时候 $x=0$ 才真是奇点？看 $y''$ 前面的系数。比如 $x^2y''+xy'+y=0$，$y''$ 的系数 $x^2$ 在 $0$ 处为零，$x=0$ 是\emph{正则奇点}，那时才要设 $y=x^r\sum a_nx^n$（Frobenius 方法），不能用本题这种普通幂级数。
}

\aside{这一句决定了后文能不能直接从 $a_0,a_1$ 起递推：常点处 $a_0,a_1$ 是两个自由常数，正好对应二阶方程的通解里的两个任意常数。}

\hdr{§2}{四个来源的系数转移}

设 $y=\sum_{n\ge0}a_nx^n$，逐项求导得
\[ y'=\sum_{n\ge0}(n+1)a_{n+1}x^{n},\qquad
   y''=\sum_{n\ge0}(n+2)(n+1)a_{n+2}x^{n}. \]

代进 $y''+(1-x)y'+(x^2-1)y$，逐项收集 $x^n$ 的系数。四个来源分别是：

\begin{itemize}
\item $y''$：$x^n$ 的系数是 $(n+2)(n+1)a_{n+2}$；
\item $y'$：$x^n$ 的系数是 $(n+1)a_{n+1}$；
\item $-xy'$：$y'$ 里 $x^{n}$ 的系数是 $(n+1)a_{n+1}$，乘 $-x$ 后落到 $x^{n+1}$，所以对 $x^n$ 的贡献是 $-n\,a_n$（$n=0$ 时为 $0$）；
\item $(x^2-1)y$：$x^2y$ 把 $a_{n-2}$ 抬到 $x^n$，$-y$ 给出 $-a_n$，合计 $a_{n-2}-a_n$（约定 $a_{-1}=a_{-2}=0$）。
\end{itemize}

把它们相加，$x^n$ 的总系数是
\[ (n+2)(n+1)a_{n+2}+(n+1)a_{n+1}-n\,a_n+a_{n-2}-a_n , \]
合并两个 $a_n$ 项，得
\[ (n+2)(n+1)a_{n+2}+(n+1)a_{n+1}-(n+1)a_n+a_{n-2}=0 . \]

\hdr{§3}{每多一项，代一次}

令上面的系数为零，就得到本题的递推：
\[
\boxed{\;(n+2)(n+1)a_{n+2}=(n+1)a_n-(n+1)a_{n+1}-a_{n-2}\;}
\qquad n=0,1,2,\dots
\]
约定 $a_{-1}=a_{-2}=0$。等价地写成“解出下一项”的形状：
\[
a_{n+2}=\frac{(n+1)a_n-(n+1)a_{n+1}-a_{n-2}}{(n+2)(n+1)} .
\]

下面这张表只是把递推式在 $n=0,\dots,7$ 上逐行读了一遍：第二列是“$x^n$ 的系数 $=0$”展开后的原方程，第三列由它解出。

{\small
\begin{tabular}{@{}r@{\quad}l@{\quad}l@{}}
\thead{$n$} & \thead{$x^n$ 的系数为零} & \thead{解出}\\ \hline
$0$ & $2a_2+a_1-a_0=0$ & $a_2=(a_0-a_1)/(2)$\\[4pt]
$1$ & $6a_3+2a_2-2a_1=0$ & $a_3=(2a_1-2a_2)/(6)$\\[4pt]
$2$ & $12a_4+3a_3-3a_2+a_0=0$ & $a_4=(3a_2-3a_3-a_0)/(12)$\\[4pt]
$3$ & $20a_5+4a_4-4a_3+a_1=0$ & $a_5=(4a_3-4a_4-a_1)/(20)$\\[4pt]
$4$ & $30a_6+5a_5-5a_4+a_2=0$ & $a_6=(5a_4-5a_5-a_2)/(30)$\\[4pt]
$5$ & $42a_7+6a_6-6a_5+a_3=0$ & $a_7=(6a_5-6a_6-a_3)/(42)$\\[4pt]
$6$ & $56a_8+7a_7-7a_6+a_4=0$ & $a_8=(7a_6-7a_7-a_4)/(56)$\\[4pt]
$7$ & $72a_9+8a_8-8a_7+a_5=0$ & $a_9=(8a_7-8a_8-a_5)/(72)$\\
\hline
\end{tabular}}

逐行读法：$n=0$ 时 $a_{-2}$ 与 $-0\cdot a_0$ 都不出现，只剩 $2a_2=a_0-a_1$；$n=1$ 时 $a_{-1}=0$，得 $6a_3=2a_1-2a_2$；从 $n=2$ 起 $a_{n-2}$ 才真正登场（$n=2$ 时是 $a_0$，$n=3$ 时是 $a_1$，依此类推）。每行的 $a_n$ 系数之所以是 $n+1$，是因为 $-n a_n$ 与 $-a_n$ 两处合并而来，这一格可以直接对照上一节。

\hdr{§4}{取 $a_0=1$,\ $a_1=0$}

取最简单的初值，把表里的式子逐行代入：
\begin{align*}
a_2&=\frac{1-0}{2}=\frac12, &
a_3&=\frac{0-1}{6}=-\frac16,\\
a_4&=\frac{3/2+1/2-1}{12}=\frac1{12}, &
a_5&=\frac{-4/6-4/12}{20}=-\frac1{20},\\
a_6&=\frac{5/12+5/20-1/2}{30}=\frac1{180}, &
a_7&=\frac{-6/20-6/180+1/6}{42}=-\frac1{252},\\
a_8&=\frac{7/180+7/252-1/12}{56}=-\frac1{3360}, &
a_9&=\frac{-8/252+8/3360+1/20}{72}=\frac{13}{45360}.
\end{align*}
于是
\[
y_1=1+\frac{x^2}{2}-\frac{x^3}{6}+\frac{x^4}{12}-\frac{x^5}{20}
+\frac{x^6}{180}-\frac{x^7}{252}-\frac{x^8}{3360}+\frac{13x^9}{45360}+\cdots
\]

\textbf{前五个非零项：}
\[ 1,\qquad \frac{x^2}{2},\qquad -\frac{x^3}{6},\qquad \frac{x^4}{12},\qquad -\frac{x^5}{20}. \]

\hdr{§5}{取 $a_0=0$,\ $a_1=1$}

同样逐行代入：
\begin{align*}
a_2&=-\frac12, &
a_3&=\frac{2-2(-1/2)}{6}=\frac12,\\
a_4&=\frac{3(-1/2)-3(1/2)-0}{12}=-\frac14, &
a_5&=\frac{4(1/2)-4(-1/4)-1}{20}=\frac1{10},\\
a_6&=\frac{5(-1/4)-5(1/10)+1/2}{30}=-\frac1{24}, &
a_7&=\frac{6/10+6/24-1/2}{42}=\frac1{120},\\
a_8&=\frac{-7/24-7/120+1/4}{56}=-\frac1{560}, &
a_9&=\frac{8/120+8/560-1/10}{72}=-\frac1{3780}.
\end{align*}
于是
\[
y_2=x-\frac{x^2}{2}+\frac{x^3}{2}-\frac{x^4}{4}+\frac{x^5}{10}
-\frac{x^6}{24}+\frac{x^7}{120}-\frac{x^8}{560}-\frac{x^9}{3780}+\cdots
\]

\textbf{前五个非零项：}
\[ x,\qquad -\frac{x^2}{2},\qquad \frac{x^3}{2},\qquad -\frac{x^4}{4},\qquad \frac{x^5}{10}. \]

\hdr{§6}{两个解的线性组合}

由取法直接读出
\[ y_1(0)=1,\quad y_1'(0)=0,\qquad y_2(0)=0,\quad y_2'(0)=1 . \]
朗斯基行列式在 $x=0$ 处为
\[ W(0)=\begin{vmatrix}1&0\\0&1\end{vmatrix}=1\neq0 , \]
故 $y_1,y_2$ 线性无关，通解为
\[ \boxed{\,y=C_1y_1+C_2y_2\,},\qquad C_1,C_2\in\mathbb{R}\ (\text{或 }\mathbb{C}). \]
等价地，$C_1=y(0)$，$C_2=y'(0)$。

\hdr{§7}{保证收敛的区间}

\callout{
\textbf{定理（常点处的幂级数解）}\;
设 $P,Q$ 在 $|x-x_0|<R$ 内解析。则方程 $y''+P(x)y'+Q(x)y=0$ 的任一解在 $|x-x_0|<R$ 内可展成收敛的幂级数 $\sum_{n\ge0}a_n(x-x_0)^n$；特别地，两个线性无关解的幂级数都在 $|x-x_0|<R$ 内收敛。
}

本题 $P=1-x$、$Q=x^2-1$ 是多项式，在全平面解析，即上式中的 $R=\infty$。所以
\[ \boxed{\text{收敛区间 }(-\infty,+\infty),\qquad \text{收敛半径 }R=\infty .} \]
直观地说：级数能走多远，由系数函数 $P,Q$ 的\textbf{奇点}决定；本题的系数函数没有奇点，级数就一路收敛到底。$y_1,y_2$ 实际上是整函数。

\aside{这条定理给的是“至少收敛到最近奇点”：$R$ 是 $P,Q$ 的最近奇点到 $x_0$ 的距离。这里 $P,Q$ 无奇点，故 $R=\infty$；不是“收敛半径由解的奇点决定”。}

\hdr{§8}{两条独立的核对}

\textbf{（1）代数验算。}
把上面 $40$ 项的部分和代回 $L[y]=y''+(1-x)y'+(x^2-1)y$，$x^0$ 到 $x^{38}$ 的系数全部为零，第一个非零系数出现在 $x^{39}$——这是截断造成的，因为递推只保证到 $x^{38}$。$y_1,y_2$ 都是如此。

\textbf{（2）数值验算。}
用高精度 ODE 求解器（DOP853，$\mathrm{rtol}=10^{-13}$、$\mathrm{atol}=10^{-16}$）解初值问题，与 $60$ 项部分和比较：

{\small
\begin{tabular}{@{}l@{\quad}r@{\quad}r@{\quad}r@{}}
\thead{解} & \thead{$x$} & \thead{级数部分和} & \thead{数值解} \\ \hline
$y_1$ & $0.7$  & $1.199757278879$ & $1.199757278879$\\
$y_1$ & $-0.5$ & $1.152720094857$ & $1.152720094857$\\
$y_1$ & $1.5$  & $1.601902644196$ & $1.601902644196$\\
$y_2$ & $0.7$  & $0.578959467963$ & $0.578959467963$\\
$y_2$ & $-0.5$ & $-0.706972224611$ & $-0.706972224611$\\
$y_2$ & $1.5$  & $1.179733176952$ & $1.179733176952$\\
\hline
\end{tabular}}

逐行读法：每行前两列是同一个初值问题的两种算法——左列把 $60$ 项级数在 $x$ 处求值，右列从 $x=0$ 积分到 $x$；两列之差都在 $10^{-13}$ 以内（表中未列出的差为 $2.9\times10^{-15}$、$-1.6\times10^{-15}$、$7.5\times10^{-14}$、$-1.1\times10^{-15}$、$1.3\times10^{-15}$、$1.2\times10^{-14}$）。$x=1.5$ 越过 $0$ 更远，差略大，与截断项数一致。

\hdr{§9}{五处最容易翻车}

\warnbox{
\textbf{1. $a_3$ 的符号。} $n=1$ 的方程是 $6a_3=2a_1-2a_2$，$a_0=1,a_1=0$ 时 $a_2=1/2$，故 $a_3=-1/6$。漏掉 $-2a_1$ 会得到 $a_3=0$。\\
\textbf{2. $n=0$ 时没有 $a_{n-2}$。} $a_{-2}=0$，别把 $a_0$ 塞进 $x^0$ 的方程；要到 $n=2$ 才出现 $a_0$。\\
\textbf{3. 别把 $x=\pm1$ 当奇点。} $x^2-1=0$ 只是 $y$ 的系数取零；奇点要看 $P,Q$ 的极点，本题没有。\\
\textbf{4. 别把 $1-x$ 看成 $1-x^2$。} 改成 $1-x^2$ 后整个递推与所有系数都会变；本题与第二题（$2x^2y''+(2x^2+x)y'+(3x-1)y=0$）是两道不同的题，别串。\\
\textbf{5. 收敛区间不是 $|x|<1$。} 本题系数函数无奇点，$R=\infty$；$|x|<1$ 是那些“系数有 $\pm1$ 奇点”的题目的答案。
}

\hdr{§10}{记住这三条就够}

\begin{enumerate}[label=\textcolor{violet}{\bfseries\arabic*.}]
\item $x=0$ 是常点（$P,Q$ 解析），所以直接设 $y=\sum a_nx^n$，$a_0,a_1$ 是自由常数。
\item 递推只牵扯 $a_{n+2}$ 与 $a_{n+1},a_n,a_{n-2}$；两组初值给两个线性无关解，通解是它们的组合。
\item 系数函数是多项式、无奇点，两个幂级数在全实轴收敛：$R=\infty$。
\end{enumerate}

\flowerrule

\subhdr{答案}
\[
y_1=1+\frac{x^2}{2}-\frac{x^3}{6}+\frac{x^4}{12}-\frac{x^5}{20}+\cdots,\qquad
y_2=x-\frac{x^2}{2}+\frac{x^3}{2}-\frac{x^4}{4}+\frac{x^5}{10}+\cdots
\]
\[
y=C_1y_1+C_2y_2,\qquad x\in(-\infty,+\infty)
\]

\subhdr{把两个容易混的问题钉死}
\begin{enumerate}[label=\textcolor{violet}{\bfseries\arabic*.}]
\item \textbf{为什么 $x=0$ 不是奇点？}\;
奇点看的是 $P,Q$ 在 $x_0$ 是否解析，而不是“某个系数在 $x_0$ 取不取零”。本题
$P(x)=1-x$、$Q(x)=x^2-1$ 在 $0$ 处解析，且 $P(0)=1\neq0$，所以 $x=0$ 是常点。
$x^2-1=0$ 的根 $\pm1$ 只是 $y$ 的系数取零，$P,Q$ 仍解析，不构成奇点。
\item \textbf{为什么收敛半径不是 $1$？}\;
$R$ 由系数函数的奇点决定，而不是由解可能出现的奇点决定。本题 $P,Q$ 是多项式，
没有有限奇点，故 $R=\infty$。只有当 $y''$ 前面的系数在 $x=0$ 也取零（例如
$x^2y''+xy'+y=0$）时，$x=0$ 才成为奇点，那时才要设 $y=x^r\sum a_nx^n$。
\end{enumerate}

\end{document}
